% 页面配置文件（中南大学实验报告模板 A4竖版有框）

% A4竖版高度为29.7cm，宽度为21.0cm

\documentclass[xiaosi]{ctexart}

% --- 页面与布局设置 ---
\usepackage[a4paper]{geometry}
\geometry{
	top=1.6cm, % 调整上边距
	bottom=1.6cm, % 调整下边距
	left=2.5cm, % 调整左边距
	right=2.0cm % 调整右边距
}
\setlength{\parindent}{2em}

% --- 字体设置 ---
\usepackage{xeCJK}
\usepackage{fontspec}
\usepackage{unicode-math}
% 设置中文字体
\setCJKmainfont{SimSun}      % 正文宋体
\setCJKsansfont{SimHei}      % 无衬线黑体
\setCJKmonofont{FangSong}    % 等宽仿宋 (可自行修改)
\setmainfont{Times New Roman} % 设置英文字体
\setmathfont{XITS Math} % 设置数学公式字体

% --- 宏包引入 ---
\usepackage{graphicx}       % 用于插入图片
\usepackage{titlesec}       % 用于自定义标题格式
\usepackage{zhnumber}       % 用于将数字转为中文序号
\usepackage{amsmath}        % AMS数学公式宏包
\usepackage{framed}         % 用于给正文加边框
\usepackage{tikz}           % 用于绘制装订线
\usepackage{eso-pic}        % 用于在页面固定位置放置内容
\usepackage{indentfirst}
\usepackage{caption}
\usepackage{setspace}
\usepackage{booktabs}
\usepackage{float}
\usepackage{listings} % listings 宏包用于代码高亮
\usepackage{xcolor} % xcolor 宏包用于设置代码颜色

% --- 页面样式与路径设置 ---
\pagestyle{empty} % 全文不显示页码
\graphicspath{{./picture/}} % 设置图片都存放在./picture/文件夹内

% --- 装订线设置 ---
% 使用eso-pic和tikz在每一页的左侧绘制虚线和文字“装订线”
\AddToShipoutPictureBG{
	\AtPageLowerLeft{
		\begin{tikzpicture}[remember picture, overlay]
			% 装订线位于左页边距3cm的中间位置
			\coordinate (start) at (1.25cm, 1.5cm); 
			\coordinate (end) at (1.25cm, 28.2cm);
			\draw[dashed] (start) -- (end);
			\node[rotate=90, anchor=south, font=\small] at (1.25cm, 14.85cm)
			{装\hspace{6em}订\hspace{6em}线};
		\end{tikzpicture}
	}
}

% --- 设置 listings 代码（代码块）样式 ---
\lstset{
	language=Python, % 设置默认语言为 Python
	basicstyle=\ttfamily\zihao{-5}, % 字体为等宽字体 (Typewriter)，五号字
	lineskip={0.5ex}, % 微调行间距，可根据效果调整
	keywordstyle=\color{blue}\bfseries, % 关键字颜色为蓝色，加粗
	commentstyle=\color[rgb]{0.0, 0.5, 0.0}\itshape, % 注释颜色为深绿，斜体
	stringstyle=\color[rgb]{0.7, 0.0, 0.0}, % 字符串颜色为深红
	showstringspaces=false, % 不显示字符串中的空格
	frame=single, % 添加单线边框
	framesep=5pt, % 边框与代码的间距
	rulesepcolor=\color{gray}, % 边框线的颜色
	backgroundcolor=\color[rgb]{0.95,0.95,0.95}, % 背景颜色为浅灰色
	breaklines=true, % 自动断行
	tabsize=4, % 制表符宽度
	numbers=left, % 显示行号
	xleftmargin=2em, % 代码框左侧留出2em空间
	xrightmargin=2em % 代码框右侧留出2em空间
}

% --- 调整公式行间距---
% 减小公式上方间距（长行）
\setlength{\abovedisplayskip}{1pt plus 0.5pt minus 0.5pt}
% 减小公式下方间距（长行）
\setlength{\belowdisplayskip}{1pt plus 0.5pt minus 0.5pt}
% 减小公式上方间距（短行）
\setlength{\abovedisplayshortskip}{1pt plus 0.5pt minus 0.5pt}
% 减小公式下方间距（短行）
\setlength{\belowdisplayshortskip}{1pt plus 0.5pt minus 0.5pt}

% --- 标题格式定义 ---
% 一级标题 (例如: 一、实验目的)
% 格式: 三号、黑体, 靠左, 缩进1个字符
\renewcommand{\thesection}{\zhnum{section}}
\titleformat{\section}
{\normalfont\heiti\zihao{3}} % 字体格式
{\thesection、}              % 标题编号
{0em}                      % 编号与标题文字的间距
{}                         % 标题前的代码
\titlespacing*{\section}
{1em}   % 左边距 (缩进1字符)
{3ex}    % 标题前间距
{2ex}    % 标题后间距

% 二级标题
% 格式: 小四、黑体, 靠左, 缩进2个字符, 需要编号
\renewcommand{\thesubsection}{（\zhnum{subsection}）}
\titleformat{\subsection}
{\normalfont\heiti\zihao{-4}}
{\thesubsection}
{0em}
{}
\titlespacing*{\subsection}{2em} {0.5ex} {0.5ex}

% 三级标题
% 格式: 小四、楷体, 靠左, 缩进2个字符, 需要编号
\titleformat{\subsubsection}
{\normalfont\kaishu\zihao{-4}}
{\thesubsubsection}
{1em}
{}
\titlespacing*{\subsubsection}{2em} {0ex} {0ex}

% --- 图表标题设置 ---
\captionsetup[figure]
{
	font={small},
	labelsep=space,
	name={图},
	position=below,
	skip=5pt
}
\captionsetup[table]
{
	font={small},
	labelsep=space,
	name={表},
	position=above,
	skip=3pt
}

% --- 多级列表设置 ---
\usepackage{enumitem} % 引入列表控制宏包
\setlist[enumerate, 1]{ % 有序列表
	topsep=0pt,
	partopsep=0pt, % 设置列表和周围文字的间距为0
	itemsep=0pt, % 设置列表项之间的间距为0
	parsep=0pt, % 设置列表项内段落之间的间距
	labelsep=0.5em,	% 调整列表项的标签与文本的间距使其看起来更自然
	leftmargin=3.5em % enumerate列表缩进3.5em
}
\setlist[itemize, 1]{ % 无序列表
	label={·},
	topsep=0pt,
	partopsep=0pt, % 设置列表和周围文字的间距为0
	itemsep=0pt, % 设置列表项之间的间距为0
	parsep=0pt, % 设置列表项内段落之间的间距
	labelsep=0em, % 调整列表项的标签与文本的间距使其看起来更自然
	leftmargin=3.5em % itemize列表缩进3.5em
}